\begin{abstract}
From CPUs to GPUs, the Roofline model has become a textbook tool for revealing a platform's performance ceiling and how well a workload matches it. Compute-in-memory (CIM), however, falls outside this picture because it violates the core Roofline assumption of separate compute and memory units. We present a resident-streaming Roofline model for CIM, in which streaming-activation and resident-weight throughputs replace compute throughput and memory bandwidth, respectively. This follows from a symmetry we observe: conventional processors pair separated hardware with homogeneous data, whereas CIM pairs unified hardware with heterogeneous data---operands residing in the array and operands streaming through it. Using NeuroSim with literature-extracted parameters, we evaluate various CIM technologies, from SRAM to emerging non-volatile memories, on this Roofline and quantify the demands of mainstream large language models across computation stages. With similar periphery, the streaming throughputs of these technologies differ by a factor of only about 10\textsuperscript{2}, whereas their resident throughputs, limited by each device's write mechanism, differ by over 10\textsuperscript{4}. The memory thus decides mainly how much reuse CIM needs, from about 2--10 input vectors per weight load on SRAM and eDRAM to about 4\texttimes10\textsuperscript{4} on NOR flash, making the reuse a workload offers an explicit criterion for choosing among technologies.
\end{abstract}
\begin{IEEEkeywords}
Compute-in-memory, roofline model, streaming intensity, data reuse, large language models.
\end{IEEEkeywords}
